\chapter{From a fusion category to a quantum chain}
\label{ch:chain}

We now return to the six-object $\mathcal H_3$ fusion category and
construct a microscopic quantum chain from it. This is the lattice
construction track introduced in the reading guide; the reduced CFTs
of the preceding chapter will serve as comparison theories.
A quantum Hamiltonian needs one further choice: which
local fusion channels should cost or release energy? This chapter makes that
choice explicit. We construct the Hilbert space of the physical Haagerup
chain, derive its local projector Hamiltonian, and explain its translation
and categorical symmetries. The required linear algebra is modest: orthogonal
projections, eigenvectors of a small adjacency matrix, and Fourier sums. The
important conceptual step is to distinguish a fusion label, a physical anyon,
and an ordinary local spin.
The fusion-space conventions of Chapter~\ref{ch:fusion} and the explicit
Haagerup associator in Chapter~\ref{ch:haagerup} are the categorical inputs.

\section{The physical object and the fusion-path basis}
\index{fusion path}\index{anyon chain}

Write the six simple objects as
\[
 I_j=\alpha^j,\qquad X_j=\alpha^j\rho,\qquad j\in\Z/3\Z.
\]
Their dimensions are $\qdim(I_j)=1$ and $\qdim(X_j)=d$, where
\[
 d^2=1+3d,\qquad d=\frac{3+\sqrt{13}}2.
\]
We put one physical object $\rho$ at every site. The physical objects are all
the same; it is their possible fusion histories that fluctuate.

An open fusion history starts with a label $a$, fuses successively with the
physical $\rho$ objects, and ends with a label $b$. At a trivalent vertex the
outgoing label must occur in the fusion of the incoming label with $\rho$.
The multiplicities here are zero or one, so no additional vertex index is
needed. For example,
\[
 I_j\rho=X_j,\qquad
 X_j\rho=I_j+X_0+X_1+X_2.
\]
An invertible height must therefore be followed by its corresponding
noninvertible height. A noninvertible height can be followed by any
noninvertible height or by its one matching invertible height.

\begin{definition}[Regular periodic path space]
The length-$L$ periodic Hilbert space $\HH_{\mathrm{path},L}$ has an orthonormal basis
\[
 |x_0,x_1,\ldots,x_{L-1}\rangle
\]
indexed by admissible closed paths, meaning that
$x_{i+1}\subset x_i\otimes\rho$ for every $i$, with indices modulo $L$.
The inner product is the ordinary orthonormal fusion-basis inner product.
\end{definition}

The word ``regular'' records that all starting heights and compatible closed
flux sectors are included. It is a choice of periodic Hilbert space, not a
synonym for the sector obtained by fusing all $L$ physical objects to the
tensor unit. Open chains with fixed endpoint labels give other spaces.

One can embed $\HH_{\mathrm{path},L}$ in $(\C^6)^{\otimes L}$ by identifying each height
with a six-state basis vector and setting the coefficients of forbidden words
to zero. This embedding is useful for ordinary Schmidt decompositions.
Chapter~\ref{ch:tensor_networks} develops those decompositions and their
tensor-network realization.
Nevertheless, $\HH_{\mathrm{path},L}$ is a constrained subspace: changing one height
arbitrarily can violate two adjacent fusion constraints. A Hamiltonian or
tensor ansatz must preserve these constraints.

\begin{example}[Counting closed paths exactly]
In the ordered height basis $(I_0,I_1,I_2;X_0,X_1,X_2)$, right fusion by
$\rho$ has adjacency matrix
\[
 N_\rho=\begin{pmatrix}0&I_3\\I_3&J_3\end{pmatrix},
\]
where $J_3$ is the all-ones matrix. Its $L$th power counts length-$L$ paths,
and taking the trace closes them. On the color-independent subspace,
$N_\rho$ is the matrix $\left(\begin{smallmatrix}0&1\\1&3\end{smallmatrix}\right)$.
Its eigenvalues are $d$ and $-d^{-1}$. On each of the two color-orthogonal
directions it is $\left(\begin{smallmatrix}0&1\\1&0\end{smallmatrix}\right)$,
with eigenvalues $1,-1$. Thus
\begin{equation}
 \dim\HH_{\mathrm{path},L}=d^L+(-d^{-1})^L+2\bigl(1+(-1)^L\bigr).
 \label{eq:path-dimension}
\end{equation}
This apparently irrational expression is an integer. Indeed, if
$t_L=d^L+(-d^{-1})^L$, then
\[
 t_0=2,\quad t_1=3,\quad t_{L+2}=3t_{L+1}+t_L.
\]
Consequently the path-space dimensions at $L=3,4,6,9$ are
$36,123,1302,46764$, respectively. The dimension grows as $d^L$,
even though the local height alphabet has six letters.
\end{example}

This calculation already explains why symmetry reduction and matrix-free
algorithms matter. A final energy table is small; the vector needed to obtain
it can have exponentially many entries. Nothing about the file size of a
reported spectrum estimates the memory needed for its calculation.

\section{An \texorpdfstring{$F$}{F}-move becomes a local projector}

Fix three consecutive heights $a,e,c$. They represent a piece of fusion tree
in which $a$ first fuses with one physical $\rho$ to $e$, and then with the
next physical $\rho$ to $c$. There is another fusion tree: first fuse the two
physical objects into $f\subset\rho^2$, and then fuse $a$ with $f$ into $c$.
The unitary matrix $F^{a\rho\rho}_c$ changes between these trees. We use
\begin{equation}
 |((a\rho)e\rho);c\rangle
  =\sum_f F^{a\rho\rho}_c[e,f]\,
       |(a(\rho\rho)f);c\rangle.
 \label{eq:chain-F-convention}
\end{equation}
The allowed $e$ and $f$ depend on the fixed endpoints. Every sum below is
over allowed channels only.

\begin{figure}[htbp]
\centering
\begin{tikzpicture}[scale=.9, every node/.style={font=\small}]
 \draw (0,0)--(.7,.8)--(1.5,1.6)--(1.5,2.35);
 \draw (1.4,0)--(.7,.8);
 \draw (2.8,0)--(1.5,1.6);
 \node[below] at (0,0) {$a$};
 \node[below] at (1.4,0) {$\rho$};
 \node[below] at (2.8,0) {$\rho$};
 \node[left] at (1.05,1.2) {$e$};
 \node[above] at (1.5,2.35) {$c$};
 \draw[-{Stealth}] (3.3,1.1)--(4.7,1.1);
 \node[above] at (4,1.1) {$F^{a\rho\rho}_c$};
 \draw (5.2,0)--(6.5,1.6)--(6.5,2.35);
 \draw (6.6,0)--(7.3,.8)--(6.5,1.6);
 \draw (8,0)--(7.3,.8);
 \node[below] at (5.2,0) {$a$};
 \node[below] at (6.6,0) {$\rho$};
 \node[below] at (8,0) {$\rho$};
 \node[right] at (6.95,1.2) {$f$};
 \node[above] at (6.5,2.35) {$c$};
\end{tikzpicture}
\caption{The local recoupling move. The neighboring heights $a,c$ are
fixed. In the left tree the middle path height $e$ varies; in the right
tree the physical pair's fusion channel $f$ is measured.}
\label{fig:chain-F-projector}
\end{figure}

In the second basis, asking whether the physical pair fuses to $\rho$ is
simply a diagonal yes-or-no measurement. Transforming that projection back
to the path basis gives
\begin{equation}
 (P^{(\rho)}_{a,c})_{e',e}
   =\overline{F^{a\rho\rho}_c[e',\rho]}
      F^{a\rho\rho}_c[e,\rho].
 \label{eq:rho-projector}
\end{equation}
If the $\rho$ channel is absent for the given endpoints, this block is zero.
Otherwise it is an outer product and has rank one.

\begin{proposition}
For each endpoint block, $P^{(\rho)}$ is an orthogonal projection. On the
periodic path space, the operator $P_i^{(\rho)}$ changes only the height
$x_i$ and preserves admissibility.
\end{proposition}
\begin{proof}
Put $v_e=\overline{F[e,\rho]}$. Unitarity gives $v^\dagger v=1$, and
\(P=vv^\dagger\). Hence $P^\dagger=P$ and
$P^2=v(v^\dagger v)v^\dagger=P$. Both the original and the updated middle
height are allowed channels with the same endpoints, so the two adjacent
fusion constraints remain satisfied. All other constraints are unchanged.
\end{proof}

The physical Hamiltonian studied here is
\begin{equation}
 H_L=-\sum_{i=0}^{L-1}P_i^{(\rho)}.
 \label{eq:physical-chain-H}
\end{equation}
The minus sign rewards the $\rho$ channel. Changing the sign, rewarding the
unit channel, or using another associator defines another dynamical problem.
Neighboring terms usually do not commute, because the fusion of sites
$i,i+1$ and that of sites $i+1,i+2$ are incompatible measurements. The
pentagon controls the consistency of all reassociations; it does not make
these overlapping measurements simultaneously diagonal.

The simple bounds $-L\leq H_L\leq0$ follow from $0\leq P_i\leq1$.
Saturating the lower bound requires a state in the $+1$ eigenspace of every
projector. Such a state need not exist at a given length. A sum of local
projections is not automatically frustration free.

\begin{example}[An exact local imaginary-time gate]
Since $P^2=P$, its exponential is
\[
 e^{\epsilon P}=1+(e^\epsilon-1)P.
\]
Thus the Trotter gate for $e^{-\epsilon H}$ on a single bond is an exact
linear combination of $1$ and $P_\rho$. Approximation enters when products
of overlapping gates replace the exponential of their sum, not in the
single-projector exponential itself. This gives a direct tensor-network
route from the Hamiltonian to a two-dimensional Euclidean model.
\end{example}

\section{Why the associator is physical data}

A change of orthonormal basis at every trivalent vertex changes the printed
$F$-symbols. If all operators and states are transformed consistently, it
does not change the spectrum. This is a gauge change. In contrast, two
inequivalent unitary associators with the same fusion rules need not give
the same local projectors or spectra.

The categories commonly called $\mathcal H_2$ and $\mathcal H_3$ provide exactly this warning.
Both have the six-object fusion ring used above. The physical chain in
\eqref{eq:physical-chain-H} uses the specified real unitary $\mathcal H_3$
associator. A uniform positive solution occurring in several algebraic
constructions belongs to $\mathcal H_2$. Substituting it because the object names
and fusion rules agree changes the chain.

\begin{constructioninput}[A finite distinction between $\mathcal H_2$ and $\mathcal H_3$]
For the same rho-channel coupling and the regular periodic length-four
path space, the repository's exact local-data calculation gives
\[
 \Tr(H_2^2)-\Tr(H_3^2)=\frac{16}{3d}>0.
\]
The space has dimension $123$, as independently obtained from
\eqref{eq:path-dimension}. This is a finite algebraic distinction between
the two Hamiltonians, not a statement about their eventual infrared limits.
\end{constructioninput}

Unitary equivalence preserves every trace of a power. The displayed
difference therefore excludes a unitary equivalence of these two finite
Hamiltonians at the stated coupling and closure. It does not exclude a
more elaborate equivalence involving module categories, boundary seams,
changed couplings, or a common infrared fixed point. Such a proposal must
specify those extra choices.

\section{Translation and the ordinary \texorpdfstring{$C_3$}{C3} charge}
\index{translation!fusion-path chain}

Let $U_{\rm tr}$ roll the height word one place to the right. Let $A$ increase every
color label by one, preserving whether the height is invertible or
noninvertible. In a compatible gauge the Hamiltonian commutes with both,
and $U_{\rm tr}^L=A^3=1$. They commute with each other. A simultaneous eigenspace
is labeled by
\[
 U_{\rm tr}=e^{2\pi im/L},\qquad A=e^{2\pi iq/3},
 \quad m\in\Z/L\Z,\quad q\in\Z/3\Z.
\]
There are $3L$ possible ordinary symmetry sectors, though their dimensions
need not be equal.

The formal sector projector is
\[
 \Pi_{m,q}=\frac1{3L}\sum_{r=0}^{L-1}\sum_{s=0}^2
       e^{-2\pi i(mr/L+qs/3)}U_{\rm tr}^rA^s.
\]
It is often implemented through orbit representatives rather than as a
dense matrix. Short orbits require care: if a combined translation and
color shift fixes a path, the chosen character must be trivial on its
stabilizer or its projected vector vanishes. Dividing indiscriminately by
$\sqrt{3L}$ gives the wrong normalization on such orbits.

For $h_j=-P_j^\rho$, our translation convention is
$U_{\rm tr}h_jU_{\rm tr}^{-1}=h_{j+1}$. Define
\begin{equation}
 H_m=\sum_{j=0}^{L-1}e^{-2\pi imj/L}h_j.
 \label{eq:energy-Fourier}
\end{equation}
Changing the summation variable gives
$U_{\rm tr}H_mU_{\rm tr}^{-1}=e^{2\pi im/L}H_m$.
Therefore $H_m$ raises momentum by $m$.
The sign can look reversed if one uses a left translation or the opposite
Fourier exponential. There is no physical disagreement, but the state
selector and Fourier operator must use one convention together.

\begin{example}[A selection rule without diagonalization]
Suppose $|a\rangle$ has charge $q_a$ and momentum $m_a$, and an operator
$O$ has charge $q_O$ and momentum $m_O$. If
$\langle b|O|a\rangle\ne0$, then
\[
 q_b=q_a+q_O\pmod3,\qquad m_b=m_a+m_O\pmod L.
\]
This follows by inserting $A^{-1}A$ and $U_{\rm tr}^{-1}U_{\rm tr}$ on both sides of the
matrix element. It is an exact lattice statement and needs no CFT limit.
\end{example}

\section{Dragging an auxiliary line around the ring}
\index{topological symmetry operator}

Ordinary color rotation is only the invertible part of the symmetry. To
construct a general categorical line, put an auxiliary object $c$ around
the periodic fusion tree and move it past each physical $\rho$ by an
$F$-move. Closing the auxiliary line gives an operator $Y_c$. In the
orthonormal path convention its matrix elements are a cyclic product of
the corresponding $F$ coefficients:
\begin{equation}
 \langle b|Y_c|a\rangle
 =\prod_i F(c,a_i,\rho,b_{i+1};b_i,a_{i+1}).
 \label{eq:topological-line-product}
\end{equation}
The six-argument notation specifies the external objects and internal
channels by
\[
 F(a,b,c,d;e,f)=F^{abc}_d[e,f].
\]
It avoids suppressing the height indices that
make the product meaningful. There is no extra quantum-dimension
prefactor in this convention.
In the selected gauge the invertible line $Y_\alpha$ is the color
rotation $A$ defined above.

Two auxiliary lines can be fused before or after being moved through the
chain. Pentagon coherence equates the two operations, giving
\[
 Y_aY_b=\sum_cN_{ab}^{c}Y_c.
\]
The same local graphical move shows that these lines commute with the
fusion-channel Hamiltonian. For self-dual $\rho$, unitarity makes
$Y_\rho$ self-adjoint. These statements use the associator, not merely
the fusion matrices considered as abstract integer matrices.

\begin{example}[Topological eigenvalues in the neutral sector]
From the fusion ring,
\[
 Y_\rho^2=1+(1+A+A^2)Y_\rho.
\]
On $A=1$ states this reduces to $Y_\rho^2=1+3Y_\rho$, so the only
possible eigenvalues are $d$ and $-d^{-1}$. The associated spectral
projections are
\[
 \Pi_+=\frac{Y_\rho+d^{-1}}{d+d^{-1}},\qquad
 \Pi_-=\frac{d-Y_\rho}{d+d^{-1}}.
\]
These formulas are useful for sector checks. They do not say that every
neutral eigenstate belongs to the vacuum conformal family. Both sectors
contain many lattice energies.
\end{example}

The operator $Y_\rho$ is not a unitary symmetry. Its eigenvalues have
different absolute values, and its fusion square is a sum of lines.
Replacing it by its polar unitary preserves a Hilbert-space operation,
but need not preserve locality or the multiplication of local fields.
One cannot infer a new ordinary CFT symmetry just by taking the polar
part of a commuting matrix.

In a conformal limit with a unique vacuum, a topological line acts on the
vacuum by its quantum dimension. The stress family is then in the
$Y_\rho=d$, $A=1$ sector because topological lines preserve the stress
tensor. Conversely, finding that an energy-density source has these
labels does not identify it as pure stress: other invariant fields can
have the same labels.
The continuum interpretation of line eigenvalues and twisted sectors
is developed in Chapter~\ref{ch:defects}; the finite fusion equations
alone do not supply that limiting dictionary.

\section{The category supplies kinematics, the Hamiltonian dynamics}

It is tempting to read a central charge directly from the categorical
dimension $d$, the number of simple objects, or the modular data of the
Drinfeld center. There is a simple obstruction to any such general rule.

\begin{proposition}[Spectator underdetermination]
Suppose a nonchiral CFT $\mathcal X$ realizes a specified Haagerup defect
action. For any independent nonchiral CFT $\mathcal B$, the product
$\mathcal X\otimes\mathcal B$ realizes the same specified action through
$D_a\otimes1$, while its central charge is $c_{\mathcal X}+c_{\mathcal B}$.
\end{proposition}
\begin{proof}
Fusion and junction maps on the first factor are unchanged when tensored
with the identity of the second. The stress tensors add, so their central
terms add. A defect-sector partition function is multiplied by the
modular-invariant spectator partition function, preserving its sector
transformation law.\end{proof}

This does not make the fixed Hamiltonian's infrared charge arbitrary.
Adding a spectator changes the microscopic system. It shows what must be
learned dynamically once the category has been chosen. The topological
central charge of a modular tensor category and the coefficient of the
nonchiral chain's Casimir energy answer different questions.

A similar distinction occurs for strange correlators. A topological bulk
tensor network supplies a state $|\Psi_{\mathcal C}\rangle$; a boundary
choice $\langle\Omega|$ supplies the amplitude
$\langle\Omega|\Psi_{\mathcal C}\rangle$. The boundary tensor contains
additional dynamics. In particular, a two-strand gate in the span
$a1+bP_1$ cannot distinguish $\rho$ from the other nonunit channels,
whereas the gate in the worked example above does. An entire decorated
strange-correlator cell also contains algebra and module tensors. Its
Hamiltonian limit must be established by an explicit local contraction,
including seams and sectors, rather than by recognizing a familiar
fusion category in one of its tensors.

\section{Exercises}

\begin{exercise}\label{ex:chain-dimension}
Use the recurrence for $t_L$ to compute $\dim\HH_{\mathrm{path},5}$ and
$\dim\HH_{\mathrm{path},8}$. Explain why the leading exponential growth is $d^L$ and
why this differs from $6^L$.

\par\smallskip\noindent{\footnotesize\normalfont
\hyperref[sol:ex:chain-dimension]{Solution on page~\pageref*{sol:ex:chain-dimension}.}}\par
\end{exercise}
\begin{exercise}\label{ex:chain-projector}
Let $F$ be a unitary $n\times n$ matrix and let $S$ be a subset of its
columns. Show that the sum of the corresponding rank-one channel
projectors is an orthogonal projector of rank $|S|$. What changes if the
fusion multiplicity of a selected channel is greater than one?

\par\smallskip\noindent{\footnotesize\normalfont
\hyperref[sol:ex:chain-projector]{Solution on page~\pageref*{sol:ex:chain-projector}.}}\par
\end{exercise}
\begin{exercise}\label{ex:chain-orbit}
For a finite abelian group $G$ acting by permutations on an orthonormal
basis, construct the character projection of one orbit vector. Prove
that it vanishes precisely when the character is nontrivial on the
orbit stabilizer, and normalize it when it is nonzero.

\par\smallskip\noindent{\footnotesize\normalfont
\hyperref[sol:ex:chain-orbit]{Solution on page~\pageref*{sol:ex:chain-orbit}.}}\par
\end{exercise}
\begin{exercise}\label{ex:chain-noninvertible}
On the alpha-neutral space prove that $Y_\rho-2$ is invertible and that
\[
 (Y_\rho-2)^\dagger(Y_\rho-2)=5-Y_\rho.
\]
Give upper and lower bounds for its squared norm on a unit vector.
Why is it incorrect to call $Y_\rho-2$ a projection on this whole space?

\par\smallskip\noindent{\footnotesize\normalfont
\hyperref[sol:ex:chain-noninvertible]{Solution on page~\pageref*{sol:ex:chain-noninvertible}.}}\par
\end{exercise}
\begin{exercise}\label{ex:chain-boundary}
Write the dimension of the open path space with fixed endpoints $a,b$
as a matrix element of $N_\rho^L$. Explain why a periodic trace and a
fixed-endpoint matrix element need not have the same low-energy spectrum.

\par\smallskip\noindent{\footnotesize\normalfont
\hyperref[sol:ex:chain-boundary]{Solution on page~\pageref*{sol:ex:chain-boundary}.}}\par
\end{exercise}
